% Diagrama hecho a mano de Comparacion_QR(); actualizar si cambia la lógica.
\begin{figure}[H]\centering
\resizebox{\linewidth}{!}{%
\begin{tikzpicture}[>=Stealth,
  st/.style={draw,rounded corners,fill=orange!10,align=center,minimum width=3.4cm,font=\footnotesize},
  dc/.style={draw,diamond,aspect=2.4,fill=yellow!15,align=center,font=\scriptsize,inner sep=1pt}]
\node[st] (read) at (0,0) {Leer tag izquierdo (id 2)\\y derecho (id 3)};
\node[dc,below=0.6cm of read] (none) {¿Ninguno\\detectado?};
\node[st,right=1.5cm of none] (cones) {Dos conos:\\\texttt{Mov\_Auto(1, 55)}};
\node[dc,below=0.6cm of none] (same) {¿Misma posici\'on\\en ambos tags?};
\node[dc,below=0.6cm of same] (cube) {¿Tag central\\= 2?};
\node[st,right=1.5cm of cube] (grab) {Agarrar cubo};
\node[st,below=0.6cm of cube] (move) {Avanzar \texttt{Mov\_Auto(1, 44)},\\releer tag y comparar};
\node[st,right=1.5cm of same] (cmp) {Comparar ambos tags\\\texttt{Comparacion\_Valores(..., 0)}};
\node[st,below=0.6cm of move] (val) {\texttt{Comparacion\_Valores(..., 1)}\\\texttt{Mov\_Auto}};
\draw[->] (read) -- (none);
\draw[->] (none) -- node[above,font=\scriptsize]{s\'i} (cones);
\draw[->] (none) -- node[right,font=\scriptsize]{no} (same);
\draw[->] (same) -- node[right,font=\scriptsize]{s\'i} (cube);
\draw[->] (same) -- node[above,font=\scriptsize]{no} (cmp);
\draw[->] (cube) -- node[above,font=\scriptsize]{s\'i} (grab);
\draw[->] (cube) -- node[right,font=\scriptsize]{no} (move);
\draw[->] (move) -- (val);
\end{tikzpicture}}
\caption{Flujo de decisi\'on de \texttt{Comparacion\_QR()} (se repite durante 28 s).}
\end{figure}
